\section{Considerações parciais e próximos passos}
\label{sec:consideracoes}

As Fases 1 e 2 transformaram percepções sobre o SIGAA/UnB em uma especificação verificável: três características da ISO/IEC 25010:2023, 16 questões com hipóteses e 16 métricas com níveis fixados antes da coleta. A Fase 3 converteu essa especificação em um plano repetível, com protocolo padronizado para o teste com usuários, inspeções por pares, regras de proteção dos dados pessoais e uma execução instrumentada que mede os critérios de Flexibilidade diretamente no navegador, sem depender da interpretação do avaliador.

Os primeiros resultados já respondem a uma das decisões que a avaliação deve apoiar, o investimento em interface responsiva: nenhuma das 12 verificações de M14 atende aos critérios, por uma causa única e estrutural (a ausência da declaração de \textit{viewport} e o \textit{layout} de 980~px), e a hipótese H15 foi confirmada. Os relatos dos dois primeiros participantes sobre o uso no celular apontam os mesmos problemas que a medição encontrou.

As limitações conhecidas são a amostra de conveniência, concentrada na FCTE e limitada a seis participantes; a verificação de 360~px por emulação, e não em celular real; a ausência, até aqui, do Microsoft Edge e do Safari em M13; os tempos de desempenho medidos de um único ponto de rede e em horário noturno; a impossibilidade de avaliar matrícula e trancamento sem alterar dados reais; e a ausência de acesso ao código-fonte.

Até a entrega final, a equipe vai: concluir as sessões P2 e P4 a P6 e as inspeções por pares (M1, M2, M9, M10 e M20); executar M13 no Edge e no Safari e M15 em um celular real; calcular todas as métricas e julgar as três características pelas regras da Fase 2, classificando as hipóteses e as relações HR1 e HR2; e selecionar a ação de melhoria pela tarefa e pela característica de pior resultado. Se o julgamento confirmar Flexibilidade como a característica mais fraca, uma candidata viável sem acesso ao código-fonte é uma extensão de navegador que declare a \textit{viewport} e adapte as páginas das tarefas a telas pequenas, cujo efeito pode ser medido pelo mesmo \textit{script} de M14. Na gestão, falta criar o marco e as \textit{issues} da Entrega 2 para medir o \textit{burndown} e rebalancear o esforço entre os membros.
